% !TEX root = _main.tex
% ========================================
% 卒業論文　本文
% ========================================

%1

\section{作品の解説}
作品タイトル：「日暮どきこそ耳を澄まして：「Echo」〜音だけの動物物語〜」

本章では，本研究で制作した音のみを用いた動物物語作品について詳述する．
具体的には，作品の構成要素である各リージョンごとに，使用した音声素材やその配置方法，および表現上の工夫について整理し，作品全体の内容を解説する．

本作品では、全149個の音響イベントを用いて物語を構成している。
本文では、物語の転換点や知覚上重要なイベントを中心に取り上げ、
個別の音響操作とその知覚的効果について論じる。
すべての音響イベントの一覧と詳細な設定については、付録\ref{app:events}に示す。

\subsection{あらすじ}
本作品は，嵐の中で墜落した一羽のハチドリが，仲間のもとへ戻ろうと奮闘する様子を描いた音のみを用いた動物物語である．

物語は，全3場面で構成されている．第1場面では，嵐の中でハチドリの群れが騒ぐ様子と，主人公であるハチドリが墜落するまでの過程が描かれる．第2場面では，地上に墜落した主人公がウサギに助けられ，巣穴へ運ばれる様子が描かれる．第3場面では，巣穴の中で雨音を聴きながら回復する主人公の様子が描かれる．

\subsection{登場キャラクター}
\tabref{tab:characters}に，本作品に登場するキャラクターの一覧を示す．
\begin{table*}[t]
  \centering
  \caption{登場キャラクター一覧}
  \label{tab:characters}
  \begin{adjustbox}{max width=\textwidth}
    \begin{tabular}{lllcll}
      \hline
        ID & キャラクター      & 種別    & 初出イベントID & 役柄                         & 初出時の状況          \\
        \hline
        C1 & ハチドリ（主人公）   & 鳥     & 8        & 物語の中心。墜落→回復→他者との関係を通じた心理変化 & 主人公中心視点に切り替わった直後の鳴き声。ドライ音質で自己存在を明示 \\
        C2 & ハチドリの群れ（仲間） & 鳥（複数） & 2--6     & 世界観・混沌状況の提示、主人公の孤立の対照      & 嵐の中で騒ぐ群れの鳴き声（俯瞰視点・リバーブ付き）          \\
        C3 & ウサギ         & 哺乳類   & 18--19   & 主人公の救助者・同行者・安定した判断主体       & 足音のみで接近（正体不明の存在として提示）              \\
        C4 & シカ          & 哺乳類   & 60       & 導き手（川への案内）、状況転換の触媒         & 接近する大型動物の足音（主体未確定）                 \\
        C5 & サル          & 哺乳類   & 117      & 誤解→理解という印象転換を担う存在          & 草の揺れによる存在予兆（正体未確定・警戒対象）            \\
        \hline
    \end{tabular}
  \end{adjustbox}
\end{table*}



